\section{Give the XOR kernel a testable native-word interface}
Implement \texttt{video\_xor\_words} and \texttt{video\_xor\_words\_msa}. Both receive
left and right word-array addresses in \texttt{\$a0} and \texttt{\$a1}, output in
\texttt{\$a2}, and a word count in \texttt{\$a3}. Counts from zero through 4096 are
valid. For a positive count, all three addresses are word-aligned and describe
readable inputs or a writable output of exactly four times the count bytes.

The output may be disjoint from both inputs or exactly equal to one of them. Inputs
may also coincide. Partial overlaps are outside this interface's caller preconditions.
Load both corresponding inputs before storing their output. This makes exact aliases
safe and lets a validated delta be applied in place. Count zero succeeds without
reading or writing either pointer, even if their values are otherwise unusable.

Success returns zero in \texttt{\$v0} and the number of nonzero XOR words in
\texttt{\$v1}. Invalid count or required alignment returns one and zero before input
access or output stores. Retain saved registers, the caller's stack pointer, every
byte outside the output window and inputs that are not aliased to output. Do not
perform file I/O inside this leaf computation.

Although valid image words have a zero high byte, this kernel computes all 32 bits.
A difference of \texttt{0x80000000} counts as one changed word. That broader numeric
contract exposes mistaken signed comparisons without changing the serialized RGB
format. When its inputs are previous and current images, the nonzero result count is the
changed-pixel count. When using the same primitive to apply a delta, that count
instead describes nonzero reconstructed words; the player need not use it.
The following six edits construct this kernel; the full file-codec assembly
construction remains a later stage of this draft.
